%-------------------------- identification ---------------------
% iheiddiss v0.0.1
% By James Hollway <james.hollway@graduateinstitute.ch>
% Based on ociamthesis v2.2 by Keith A. Gillow <gillow@maths.ox.ac.uk>

% Setup document % % %
\documentclass[12pt,a4paper]{article}
\usepackage{geometry}
\geometry{dvips,a4paper,margin=0.75in}
\usepackage{ragged2e} % this offers options for centering, justification, etc.

% Setup fonts % % %
\usepackage[T1]{fontenc}
%\usepackage[utf8]{inputenc}
\usepackage{microtype} % this makes fonts almost imperceptibly smoother
\usepackage{fontspec}
% For the headings we will use Helvetica
\usepackage{helvet}
\usepackage{sectsty}
\allsectionsfont{\sffamily}
% For the main text and equations we will use Baskerville and Palatino
\usepackage{mathpazo}
\usepackage{baskervald}
% For formatting code or package names, we will use Lucida Console
\usepackage{zi4}

% Setup figures % % %
\usepackage[]{graphicx}

% Setup bibliography % % %
\usepackage[backend=biber, % if this doesn't work for you, use bibtex8 as a backend
%style=verbose, % use this if you are in the law department
style=authoryear, % use this if you are in any other department
% change the following only if you know what you are doing:
maxcitenames=3, % how many authors to cite before resorting to ``et al''
maxbibnames=99, % in the bibliography we want them all though
sortcites=true, % this re-sorts citations when you cite several at once
hyperref=true, % this adds a link to a citations date that points to the reference
%backref=true, % references state on which page they are cited
abbreviate=true,
url=false, % don't add (lengthy) URL information
doi=false, % don't add (lengthy) DOI information
autolang=hyphen]{biblatex}
% The following is where you put the path to your bibliography:
\addbibresource{~/Dropbox/library.bib}

% Setup links % % %
% Keep the following in the following order, or else they may break!
\usepackage{varioref}
\usepackage[colorlinks=true,citecolor=red,linkcolor=red,urlcolor=red]{hyperref} % Let's use IHEID red
\usepackage[hypcap]{caption}
\usepackage{cleveref}

% % % % % % % % % % % % %
% Document Starts Here % % %
% % % % % % % % % % % % %
\begin{document}

% Title page
\hspace*{-1.25cm}
\includegraphics[width=0.4\linewidth]{Logo_CMYK_Hi.eps}

%\maketitle
{\centering \sffamily
\vspace{3cm}
{\Large\textbf{TITLE OF THESIS}}

\vspace{3cm}
Name Surname

\vspace{8cm}
Submitted in partial fulfilment of the requirements for the\\
	% Pick one of the following:
%	LLM in International Law\\
%	MA in International Relations/Political Science\\
%	Master in Anthropology and Sociology\\
%	Master in Development Studies\\
%	Master in International Affairs\\
%	Master in International Economics\\
%	Master in International History\\
%	Master in International Law\\
%	PhD in Anthropology and Sociology\\
%	PhD in Development Economics\\
%	PhD in International Economics\\
%	PhD in International History\\
%	PhD in International Law\\
	PhD in International Relations/Political Science\\

\vspace{1cm}Geneva

\vspace{1cm}\the\year

}
\pagenumbering{gobble}

\newpage
\pagenumbering{arabic}
\justifying

\section{Intro}

\textcite{Hollway:2017ke} is an example of an in-text citation. If you use ``autocite'' \autocite{Hollway:2017ke}, then the style will automatically update when you change it's configuration in the \TeX~ file. To learn more about how to use ``biblatex'', the package we are using here for referencing, see \href{http://tug.ctan.org/info/biblatex-cheatsheet/biblatex-cheatsheet.pdf}{this cheatsheet}.

If you want to change fonts, check \href{https://www.tug.org/pracjourn/2006-1/schmidt/schmidt.pdf}{here}.

Math can be inline such as $ y = \alpha + \beta x + \epsilon $ or look like this: $$ \bar{x} = \frac{1}{n} \sum_{i=1}^{n}x_i $$

The math might express an equation used in a \texttt{package}. Indeed, you can write whole sections of code like so:

\texttt{model1 <- estimate(y \textasciitilde~ x1 + x2, data)}
	
\newpage
\printbibliography

\end{document}